\section{Discussion and Outlook}
\label{sec:discussion}

We have shown that the common strategy of treating reduced algebraic objects (instantaneous
states, channels, or static multi-time tensors) as self-sufficient descriptions is structurally
limited. When the reduction map is degenerate in a way that discards causally relevant
correlations, and when the intervention class is rich enough to probe that dependence, the
resulting $\Sigma_1$ description is unfaithful: it cannot support counterfactual prediction
under any evolution rule defined solely on the reduced object. This limitation is
representational rather than dynamical and does not imply that $\Sigma_1$ descriptions are
universally invalid.

We then introduced $\Sigma_2$ geometry as an explicit representation of the minimal information
any faithful reduced description must encode: a base point $r$, its fiber
$\mathcal{F}(r)=\Pi^{-1}(r)$, and an intervention-indexed continuation structure
$\{\mathcal{A}_I(r)\}$, together with a set-valued update rule induced by interventions.
This construction restores counterfactual semantics by design and establishes minimality in a
\emph{representational} sense: any faithful completion must carry at least enough information to
reconstruct the $\Sigma_2$ continuation structure.

\subsection{What we are not claiming}

\begin{itemize}
  \item We do not claim that quantum theory is incorrect or incomplete as a predictive framework.
  \item We do not introduce new microscopic dynamics or modification rules.
  \item We do not propose hidden variables as an ontological commitment.
  \item We do not claim that $\Sigma_2$ descriptions are always computationally tractable in full generality.
\end{itemize}

\subsection{Open questions}

Several directions merit future work:
\begin{enumerate}
  \item \textbf{Categorical and sheaf-theoretic formulations.} The fiber/admissibility structure
        suggests natural connections to fibrations and sheaf-like semantics for local predictive
        patches.
  \item \textbf{Complexity of $\Sigma_2$ approximation.} Characterizing when low-dimensional or
        efficiently sampleable approximations exist.
  \item \textbf{Experimental diagnostics.} Designing continuation-based tests that isolate
        unfaithful cuts in controlled settings such as open quantum systems, quantum control,
        and contextuality experiments.
  \item \textbf{Intervention-class dependence.} Systematically classifying which intervention
        classes render a given reduction faithful or unfaithful.
\end{enumerate}

\subsection{A practical message}

When $\Sigma_1$ methods succeed, they do so in restricted regimes where discarded correlations
are dynamically irrelevant for the intervention class of interest. When they fail, the correct
response is not merely to search for a better reduced evolution equation, but to recognize that
the operational reduction has created a blind spot: counterfactual prediction requires tracking,
at minimum, fiber-level admissibility structure. $\Sigma_2$ provides a principled diagnostic
language for identifying this boundary and for guiding the construction of approximations that
respect it.